\documentclass[10pt,a4paper]{amsart}
\usepackage[margin=19mm]{geometry}
\usepackage{amsmath,amssymb,mathtools}
\usepackage[scr=rsfs,cal=euler]{mathalfa}
\usepackage{xfrac}
\usepackage[unicode,hidelinks,bookmarks=false]{hyperref}
\newcommand{\ie}{i.e.\ }
\newcommand{\cf}{cf.\ }
\newcommand{\resp}{respectively\ }
\newcommand{\Subsection}{\subsection}
\providecommand{\nobreakdash}{\nobreak\hbox{-}}
\setlength{\emergencystretch}{2em}
\multlinegap0pt
\usepackage{tikz}
\usepackage{amssymb}
\usepackage{bm}
\usepackage{float}
\usepackage{multirow}
\usepackage{caption}
\usepackage[shortlabels]{enumitem}
\newtheorem{theorem}{Theorem}
\newcommand{\mF}{\mathcal{F}}
\title{\textbf{Knot Theory and Error-Correcting Codes}}
\author{Altan B. K\i l\i\c{c}, Anne Nijsten, Ruud Pellikaan, Alberto Ravagnani}
\date{Source: arXiv:2307.14882v2 (CC BY 4.0)}
\begin{document}
\maketitle

\noindent\emph{Source and licence.} \textbf{Knot Theory and Error-Correcting Codes}. Version arXiv:2307.14882v2 (CC BY 4.0), distributed by arXiv under the Creative Commons Attribution 4.0 International licence, http://creativecommons.org/licenses/by/4.0/, as recorded in the arXiv licence metadata for this exact version and shown on the abstract page https://arxiv.org/abs/2307.14882v2. Full licence notice: \texttt{licenses/arxiv-2307.14882-CC-BY-4.0.txt}.\par\smallskip

\noindent\textit{Editorial scope.} Counted unit: the article's theorem thm:dimreidemeister (source0.tex lines 1372--1376). The article's complete original proof of it is delivered with it (source0.tex lines 1378--1402, one \texttt{proof} environment of 25 lines). No mathematics is written, repaired or re-derived: the statement and the proof are the article's own lines, in the article's own notation, and no retained line is substituted or re-flowed.  Retained uncounted as the source-local context the proof needs: lines 449--451; lines 1364--1369.  One declared adaptation: exactly 1 ancillary strings inside the retained spans are omitted (image-inclusion commands and, where the source repeats a label, the repeated label definitions) and no other line differs from the source; the figure environments, with their placement, captions and labels, are kept, and the package records the omitted strings in fidelity receipts bound to the affected bodies.  No retained line contains a citation, so the wrapper carries no bibliography and no bibliography entry.  The retained lines contain the article's own diagram code, which the wrapper typesets with the standard tikz packages; no image file is delivered and the package declares no asset dependency.  Editorial formatting only, in addition: an AMS \texttt{amsart} preamble that restates the article's own declarations and macros used by the retained lines, namely the environments \texttt{theorem} on separate counters, together with the standard packages those lines need.  The retained environments print in this document as Theorem 1 in order of appearance, whereas the article prints the same items with the numbering of its own full document.  The complete native source of the article as distributed in the arXiv e-print for this exact version is bundled unchanged as \texttt{sources/144859/source0.tex}.
\begin{theorem}\label{thm:reidemeistermoves}
Let $D$ and $D'$ be the diagrams of two knots $K$ and $K'$, respectively. Then~$K \approx K'$ if and only if $D \approx D'$.
\end{theorem}

\begin{figure}[ht]
    \centering
    
    \caption{The effect of the Reidemeister moves on a Fox-coloring.}
    \label{fig:fox_reidemeister}
\end{figure}

\begin{theorem}
\label{thm:dimreidemeister}
Let $D$ and $D'$ be equivalent knot diagrams. Then 
$\mF_{D,t}$ and $\mF_{D',t}$ have the same dimension.
\end{theorem}

\begin{proof}
Let $D$ and $D'$ have $n$ strands. Denote by $\mF_{D,t}$ and $\mF_{D',t}$ the respective codes related to the diagrams. One locally investigate what happens when performing each Reidemeister move, see 
Figure~\ref{fig:fox_reidemeister}. Suppose $D'$ is obtained from $D$ by twisting a strand $x_1$ (Reidemeister move of type I), then the twist results in two strands and a crossing in this part of the diagram, where both strands are the understrands and one of the strands is the overstrand. For a Fox coloring it then follows that the colors assigned to both strands must be the same. 
    Let 
    $$G = 
    \begin{pmatrix}
     | & | & | \\
    \dots & x_1 & \dots \\
    | & | & |
    \end{pmatrix}
    $$ 
    be a full rank $k \times n$ generator matrix of $\mF_{D,t}$.
    Then    $$G' = 
    \begin{pmatrix}
    | & | & | & |\\
    \dots & x_1 & x_1 & \dots\\
    | & | & | & | 
    \end{pmatrix}
    $$ 
    is a $k \times (n+1)$ generator matrix for $\mF_{D'}$, which has the same rank since the added column is a duplicate of another column.

    The other moves can be investigated
    in a similar manner and we omit the proof here. By Theorem \ref{thm:reidemeistermoves},
any two diagrams of a knot can be transformed into each other using Reidemeister moves. It follows that $\mF_{D,t}$ and $\mF_{D',t}$ have the same dimension.
\end{proof}

\end{document}
